% ============================================================
% §2 数据源溯源与精细化特征工程
% 竞赛: 电工杯 B题 (2026)
% 阶段: Stage 01 产物, L1 rubric 待审
% ============================================================

\section{数据源溯源与精细化特征工程}

\subsection{数据资产全景}

原始数据含附件1--5，覆盖某城市10个连片小区（A--J）的5类异构数据：常住人口31~300人，60+老龄人口6~864人，老龄化率21.9\%。数据结构为$10\times 3\times 6$三维需求张量——小区$\times$老人类型$\times$服务项目，5个Excel工作簿、9张表、210条记录，缺失率0.0\%。各附件谱系与统计摘要见表\ref{tab:data_inventory}。

\begin{table}[htbp]
\centering
\footnotesize
\caption{附件数据资产清单与描述性统计摘要}
\label{tab:data_inventory}
\begin{tabular}{p{1.0cm}p{3.0cm}p{1.3cm}p{5.0cm}}
\toprule
\textbf{附件} & \textbf{数据表} & \textbf{维度} & \textbf{关键统计量} \\
\midrule
附件1 & 人口与老人结构 & $10 \times 7$ & 老龄化率21.9\%, 人均月收入均值3250元 \\
       & 转移概率 & $2 \times 1$ & 自理$\to$半失能4.5\%, 半失能$\to$失能10\% \\
\midrule
附件2 & 月均服务需求次数 & $6 \times 3$ & 助餐需求最大, 紧急救助最小(0.15--3次/月) \\
       & 服务营收与支出 & $6 \times 2$ & 上门护理单价最高(30元), 紧急救助免费 \\
       & 月服务消费上限 & $3 \times 1$ & 自理$\leq$20\%, 半失能$\leq$25\%, 失能$\leq$30\% \\
\midrule
附件3 & 建设与运营成本 & $3 \times 3$ & 小/中/大型建设18/32/45万, 日管理2000/3200/4400元 \\
\midrule
附件4 & 小区间距离矩阵 & $10 \times 10$ & IQR [300,1800]m, 中位数700m, 77.8\%边$\leq$1000m \\
\midrule
附件5 & 满意度评分规则 & 3因素$\times$4-5档 & $S=0.2S_1+0.3S_2+0.5S_3$, 值域$[0.6,1.0]$ \\
\bottomrule
\end{tabular}
\end{table}

\subsection{人口结构空间异质性}

10个小区老人结构非均匀：C区920人居首，F区472人居末，极差448人。自理老人占68.6\%（4~712人）、半失能22.3\%（1~528人）、失能9.1\%（624人），与《中国老龄事业发展报告》城市社区宏观分布——自理65--70\%——高度吻合。半失能与失能合计31.3\%——近三分之一老年人口对上门护理、康复理疗等高频刚需存在强制性依赖。

\subsection{服务需求的经济约束预诊断}

自理老人月均全量消费356元，半失能822元，失能1~208元——后者的消费强度为自理老人的3.4倍。以人均月收入3~250元为基准进行消费约束预诊断：自理老人月预算650元$>$356元，充足率54.8\%，约束未触发；半失能老人月预算813元$\approx$822元，偏差仅1.1\%，F、D等低收小区已触发削减；失能老人月预算975元$<$1~208元，缺口23.9\%，\textbf{全10小区全部触发}，F区缺口高达49.1\%。此”因贫失能”张力——失能越深、需求越大、预算越紧——构成Q1.3与Q3的核心经济学驱动。

\subsection{空间距离矩阵可达性检验}

距离矩阵$10\times 10$对称，4处轻微三角不等式违反——幅度$<$200m，4/360路径——保留为道路网络迂回效应。以1000m为阈值，90条边中70条可达，占比77.8\%；A和F各仅4个可达邻居，为全区域最低；C/D/G/H/I/J各8个，为最高。此异质性直接影响Q2候选站点优先序：部署于C或J区可最大化辐射范围。

\subsection{满意度规则离散化特征}

$S = 0.2S_1 + 0.3S_2 + 0.5S_3$，其中$S_1$为4档、$S_2$为5档、$S_3$为4档阶梯分段常数，值域$[0.6,1.0]$。离散阶梯要求§3 Big-M精确线性化，不可简化为连续近似。

\subsection{预处理结论}

5附件9表210条记录零缺失、三类老人数交叉校验通过；距离矩阵对称健壮；\textbf{失能全10小区触发等比例削减}、半失能2小区触发；满意度离散阶梯确认须Big-M编码。
